\documentclass{article}
\usepackage{graphicx} % Required for inserting images

\title{Lab1 Deep Learning: Gradient Descent}
\author{Nguyen Thanh Hung}
\date{May 2026}

\begin{document}

\maketitle

\section{How you implement the algorithm?}

First, I define the $f(x) = x^2$ and $f'(x) = 2x$. 
Then, I define the gradient descent function, with x = x - lr * f'(x). 

\section{Analyze the effect of different learning rate r}

I ran with the learning rate from 0.1 to 0.9, with step 0.1. x = 10, run for 10 iterations. 

\begin{itemize}
    \item lr = 0.1: $f(x)$ decreases smoothly. It requires many steps to reach $f(x) \approx 0.013$ as it approaches the minimum from one side.
    \item lr = 0.2: Faster convergence, reaching $f(x) \approx 0.0036$ before stopping.
    \item lr = 0.3: Reaches $f(x) \approx 1.09 \times 10^{-6}$ very quickly.
    \item lr = 0.4: Stops almost immediately at $f(x) \approx 1.05 \times 10^{-12}$.
    \item lr = 0.5: Perfect convergence. Reaches $f(x) = 0$ in exactly one step.
    \item lr = 0.6: Reaches $f(x) \approx 1.05 \times 10^{-12}$ quickly, but $x$ alternates signs (zigzagging).
    \item lr = 0.7: Reaches $f(x) \approx 4.09 \times 10^{-5}$ with noticeable oscillation.
    \item lr = 0.8: $f(x)$ oscillates heavily, stopping at $f(x) \approx 0.0036$ after jumping across the minimum multiple times.
    \item lr = 0.9: Slowest "oscillating" convergence. 
\end{itemize}

\end{document}
